\documentclass[../../../main.tex]{subfiles}

\begin{document}

\section{Gaussian Process}\label{Appendix:GP}

The problem of machine learning (ML) is induction: We need to move from a finite training data $\mathcal{D}$ to a function $f$ that makes predictions for all possible input values. A common method is to restrict the class of functions considered, for example only considering linear functions of the input. Another common method is to give a prior probability to all possible functions, where higher probabilities are given to functions that we consider to be more likely, for example because they are smoother than other functions. The Gaussian process (GP) belongs to the latter approach \autocite{Rasmussen2005}.

A Gaussian process is a generalization of the Gaussian probability distribution. Whereas a probability distribution describes random variables which are scalars or vectors, a stochastic process governs the properties of functions. A function can be loosely thought of as a very long vector, each entry in the vector specifying the function value $f(x)$ at a particular input $x$. It turns out that if we only ask for the properties of the function at a finite number of points, then inference in the Gaussian process will give us the same answer as if we had considered all points \autocite{Rasmussen2005}.

% Gaussian process
\begin{figure}[!htbp]
\centering
\includegraphics[width=0.6\linewidth]{appendix/methods/pipelines/assets/gp/gp.png}
\caption[Gaussian process]{Example of a Gaussian process. Data points (red) are used to fit estimator functions (cyan). We see a distribution of functions \autocite{Fonnesbeck2022}.}
\end{figure}

\subsection{Preliminaries}

There are two equivalent ways to interpret Gaussian processes for models. One can think of a Gaussian process as defining a distribution over functions, and inference taking place in function space: the function space view. Alternatively, the weight space view considers distributions over the weights of models. Let us consider standard linear regression as an example \autocite{Rasmussen2005}.

\subsubsection{Weight Space View}

Say we have a training set $\mathcal{D}$ of $n$ observations, $\mathcal{D} = \{ (\bm{x}_i,y_i) \mid i=1,\ldots,n\}$, where $\bm{x}$ denotes an input vector of dimension $D$ and $y$ denotes a scalar output or target. The column vector inputs for all $n$ cases are aggregated in the $D \times n$ design matrix $\bm{X}$, and the target are collected in the vector $\bm{y}$, so we can write $\mathcal{D}=(\bm{X},\bm{y})$. The standard linear regression model with Gaussian noise is written as follows where $\bm{x}$ is the input vector, $\bm{w}$ is a vector of weights (parameters) of the linear model, $f$ is the function value, and $y$ is the observed target value \autocite{Rasmussen2005}.

\begin{equation}
f(\bm{x}) = \bm{x}^\top\bm{w}, \quad y = f(\bm{x}) + \epsilon.
\end{equation}

Often a bias weight or offset is included, but this can be implemented by augmenting the input vector $\bm{x}$ and weight vector $\bm{w}$, and so is not explicitly included in this notation. We have assumed that the observed values $y$  differ from the function values $f(\bm{x})$ by additive noise, and further assumed that this noise follows and independent and identically distributed (IID) Gaussian distribution with zero mean and variance $\sigma_n^2$ \autocite{Rasmussen2005}.

\begin{equation}
\epsilon \sim \mathcal{N}(0, \sigma_n^2).
\end{equation}

This gives rise to the likelihood, the probability density of the observations given the parameters, which is factored as follows by the independence assumption \autocite{Rasmussen2005}.

\begin{equation}
\begin{split}
p(\bm{y}|\bm{X},\bm{w}) &= \prod_{i=1}^{n}p(y_i|\bm{x}_i,\bm{w}) \\
&= \prod_{i=1}^{n} \frac{1}{\sqrt{2\pi}\sigma_n} \exp\left(-\frac{(y_i-\bm{x}_i^\top\bm{w})^2}{2\sigma_n^2}\right) \\
&= \frac{1}{(2\pi\sigma_n^2)^{n/2}} \exp\left(-\frac{1}{2\sigma_n^2}\left|\bm{y}-\bm{X}^\top\bm{w}\right|^2\right) \\
&= \mathcal{N}(\bm{X}^\top\bm{w},\sigma_n^2\bm{I}).
\end{split}
\end{equation}

For Bayesian formalism we specify a prior over the parameters using a zero mean Gaussian with covariance matrix $\bm{\Sigma}_p$ on the weights \autocite{Rasmussen2005}.

\begin{equation}
\bm{w} \sim \mathcal{N}(\bm{0},\bm{\Sigma}_p).
\end{equation}

Inference in the Bayesian linear model is based on the posterior distribution over the weights computed by Bayes' rule \autocite{Rasmussen2005}.

\[
\text{posterior} = \frac{\text{likelihood}\times\text{prior}}{\text{marginal likelihood}}, \quad p(\bm{w}|\bm{y},\bm{X}) = \frac{p(\bm{y}|\bm{X},\bm{w})p(\bm{w})}{p(\bm{y}|\bm{X})}.
\]

The normalizing constant known as the marginal likelihood or evidence is independent of the weights and given by the following \autocite{Rasmussen2005}.

\begin{equation}
p(\bm{y}|\bm{X}) = \int p(\bm{y}|\bm{X},\bm{w})p(\bm{p})\mathrm{d}\bm{w}.
\end{equation}

The posterior combines the likelihood and the prior capturing everything we know about the parameters, and writing only the terms from the likelihood and prior which depend on the weights we obtain the following where $\bm{\bar{w}}=\sigma_n^{-2}(\sigma_n^{-2}\bm{X}\bm{X}^\top+\bm{\Sigma}_p^{-1})^{-1}\bm{X}\bm{y}$ \autocite{Rasmussen2005}.

\begin{equation}
\begin{split}
p(\bm{w}|\bm{X},\bm{y}) &\propto \exp\left(-\frac{1}{2\sigma_n^2}(\bm{y}-\bm{X}^\top\bm{w})^\top(\bm{y}-\bm{X}^\top\bm{w})\right) \exp\left(-\frac{1}{2}\bm{w}^\top\bm{\Sigma}_p^{-1}\bm{w}\right) \\
&\propto \exp\left(-\frac{1}{2}(\bm{w}-\bm{\bar{w}})^\top(\frac{1}{\sigma_n^2}\bm{X}\bm{X}^\top+\bm{\Sigma}_p^{-1})(\bm{w}-\bm{\bar{w}})\right).
\end{split}
\end{equation}

The form of the posterior distribution is also Gaussian with mean $\bm{\bar{w}}$ and covariance matrix $\bm{A}^{-1}$ where $\bm{A}=\sigma_n^{-2}\bm{X}\bm{X}^\top+\bm{\Sigma}_p^{-1}$ \autocite{Rasmussen2005}.

\begin{equation}
p(\bm{w}|\bm{X},\bm{y}) \sim \mathcal{N}\left(\frac{1}{\sigma_n^2}\bm{A}^{-1}\bm{X}\bm{y},\bm{A}^{-1}\right).
\end{equation}

To make predictions for a test case we average over all possible parameter values weighted by their posterior probability. This is in contrast to non-Bayesian schemes where a single parameter is typically chosen by some optimization criterion. Thus the predictive distribution for $f_\ast=f(\bm{x}_\ast)$ at $\bm{x}_\ast$ is given by averaging the output of all output linear models with respect to the Gaussian posterior. The predictive distribution is again Gaussian \autocite{Rasmussen2005}.

\begin{equation}
\begin{split}
p(f_\ast|\bm{x}_\ast,\bm{X},\bm{y}) &= \int p(f_\ast|\bm{x}_\ast,\bm{w})p(\bm{w}|\bm{X},\bm{y})\mathrm{d}\bm{w} \\
&= \int \bm{x}_\ast^\top\bm{w}p(\bm{w}|\bm{X},\bm{y})\mathrm{d}\bm{w} \\
&= \mathcal{N}\left(\frac{1}{\sigma_n^2}\bm{x}_\ast^\top\bm{A}^{-1}\bm{X}\bm{y},\bm{x}_\ast^\top\bm{A}^{-1}\bm{x}_\ast\right).
\end{split}
\end{equation}

The Bayesian linear model suffers from limited expressiveness. A simple idea to overcome this problem is to first project the inputs into some high dimensional space using a set of basis functions and then apply the linear model in this space instead of directly on the inputs themselves. For example, a scalar input $x$ could be projected into the space of powers of $x$: $\phi(x)=(1,x,x^2,x^3,\ldots)^\top$ to implement polynomial regression. This idea is also used in classification where a dataset which is not linearly separable in the original data space may become linearly separable in a high dimensional feature space \autocite{Rasmussen2005}.

We assume a given basis function $\phi(\bm{x})$ which maps a $D$ dimensional input vector $\bm{x}$ into an $N$ dimensional feature space. Further let the matrix $\Phi(\bm{X})$ be the aggregation of columns $\phi(\bm{x})$ for all cases in the training set. Now the model is as follows where the vector of parameters now has length $N$ \autocite{Rasmussen2005}.

\begin{equation}
f(\bm{x}) = \phi(\bm{x})^\top\bm{w}.
\end{equation}

The predictive distribution becomes the following with $\Phi=\Phi(\bm{X})$ and $\bm{A}=\sigma_n^{-2}\Phi\Phi^\top+\bm{\Sigma}_p^{-1}$ \autocite{Rasmussen2005}.

\begin{equation}
f_\ast|\bm{x}_\ast,\bm{X},\bm{y} \sim \mathcal{N}\left(\frac{1}{\sigma_n^2}\phi(\bm{x}_\ast)^\top\bm{A}^{-1}\Phi\bm{y},\phi(\bm{x}_\ast)^\top\bm{A}^{-1}\phi(\bm{x}_\ast)\right).
\label{Equation:FeatureSpace}
\end{equation}

\autoref{Equation:FeatureSpace} can be rewritten in the following way where $\phi(\bm{x}_\ast)=\phi_\ast$ and $\bm{K}=\Phi^\top\bm{\Sigma}_p\Phi$ \autocite{Rasmussen2005}.

\begin{equation}
\begin{split}
f_\ast|\bm{x}_\ast,\bm{X},\bm{y} \sim \mathcal{N}(&\phi_\ast^\top\bm{\Sigma}_p\Phi(\bm{K}+\sigma_n^2\bm{I})^{-1}\bm{y},\\
&\phi_\ast^\top\bm{\Sigma}_p\phi_\ast-\phi_\ast^\top\bm{\Sigma}_p\Phi(\bm{K}+\sigma_n^2\bm{I})^{-1}\Phi^\top\bm{\Sigma}_p\phi_\ast).
\end{split}
\label{Equation:AltFeatureSpace}
\end{equation}

To make predictions using \autoref{Equation:FeatureSpace} we need to invert the $\bm{A}$ matrix of size $N \times N$ which may not be convenient if the dimension of the feature space $N$ is large. In \autoref{Equation:AltFeatureSpace} we need to invert matrices of size $n \times n$ which is more convenient when $n < N$. The feature space in \autoref{Equation:AltFeatureSpace} only appears in the form of $\Phi^\top\bm{\Sigma}_p\Phi$, $\phi_\ast^\top\bm{\Sigma}_p\Phi$, or $\phi_\ast^\top\bm{\Sigma}_p\phi_\ast$. Thus the entries of these matrices are invariably of the form $\phi(\bm{x})^\top\bm{\Sigma}_p\phi(\bm{x}')$ where $\bm{x}$ and $\bm{x}'$ are in either the training of the test sets \autocite{Rasmussen2005}.

Let us define $k(\bm{x},\bm{x}') = \phi(\bm{x})^\top\bm{\Sigma}_p\phi(\bm{x}')$. We call $k(\bm{x},\bm{x}')$ a covariance function or kernel. Notice that $\phi(\bm{x})^\top\bm{\Sigma}_p\phi(\bm{x}')$ is an inner product with respect to $\bm{\Sigma}_p$. As $\bm{\Sigma}_p$ is positive definite we can define $\bm{\Sigma}_p^{1/2}$. Then defining $\psi(\bm{x})=\bm{\Sigma}_p^{1/2}\phi(\bm{x})$ we obtain a simple dot product representation $k(\bm{x},\bm{x}')=\psi(\bm{x})\psi(\bm{x}')$ \autocite{Rasmussen2005}.

If an algorithm is defined solely in terms of inner products in input space then it can be lifted into feature space by replacing occurrences of those inner products by $k(\bm{x},\bm{x}')$. This is called the kernel trick. This technique is particularly valuable in situations where it is more convenient to compute the kernel than the feature vectors themselves. This often leads to considering the kernel as the object of primary interest, and its corresponding feature space as having secondary practical importance \autocite{Rasmussen2005}.

\subsubsection{Function Space View}

An alternative and equivalent way of reaching identical results to the weight space view is possible by considering inference directly in function space, and we can use a Gaussian process to describe a distribution over functions \autocite{Rasmussen2005}.

\begin{definition}
A Gaussian process is a collection of random variables, any finite number of which have a joint Gaussian distribution \autocite{Rasmussen2005}.
\end{definition}

A Gaussian process is completely specified by its mean function and covariance function, which we define as $m(\bm{x})$ and $k(\bm{x},\bm{x}')$ respectively for a real process $f(\bm{x})$ \autocite{Rasmussen2005}.

\begin{equation}
\begin{split}
m(\bm{x}) &= \mathbb{E}[f(\bm{x})] \\
k(\bm{x},\bm{x}') &= \mathbb{E}[(f(\bm{x})-m(\bm{x}))(f(\bm{x}')-m(\bm{x}'))].
\end{split}
\end{equation}

We write a Gaussian process as follows \autocite{Rasmussen2005}.

\begin{equation}
f(\bm{x}) \sim \mathcal{GP}(m(\bm{x}),k(\bm{x},\bm{x}')).
\end{equation}

A simple example of a Gaussian process can be obtained from our Bayesian linear regression model $f(\bm{x})=\phi(\bm{x})^\top\bm{w}$ with prior $\bm{w}\sim\mathcal{N}(\bm{0},\bm{\Sigma}_p)$ \autocite{Rasmussen2005}.

\begin{equation}
\begin{split}
\mathbb{E}[f(\bm{x})] &= \phi(\bm{x})^\top\mathbb{E}[\bm{w}] = 0 \\
\mathbb{E}[f(\bm{x})f(\bm{x}')] &= \phi(\bm{x})^\top\mathbb{E}[\bm{w}\bm{w}^\top]\phi(\bm{x}') = \phi(\bm{x})^\top\bm{\Sigma}_p\phi(\bm{x}').
\end{split}
\end{equation}

Thus $f(\bm{x})$ and $f(\bm{x}')$ are jointly Gaussian with zero mean and covariance given by $\phi(\bm{x})^\top\bm{\Sigma}_p\phi(\bm{x}')$. Our running example of a covariance function will be the squared exponential (SE) covariance function. For this particular covariance function, we see that the covariance is almost unity between variables whose corresponding inputs are close, and decreases as their distance in the input space increases \autocite{Rasmussen2005}.

\begin{equation}
\cov(f(\bm{x}_p),f(\bm{x}_1)) = k(\bm{x}_p,\bm{x}_q) = \exp\left(-\frac{1}{2}\left|\bm{x}_p-\bm{x}_q\right|^2\right).
\label{Equation:ExpQuad}
\end{equation}

The specification of the covariance function implies a distribution over functions. We can draw samples from the distribution of functions evaluated at any number of points. We choose a number of input points $\bm{X}_\ast$ and write the corresponding covariance matrix elementwise using \autoref{Equation:ExpQuad}. Then we generate a random Gaussian vector with this covariance matrix and further can plot the generated values as a function of the inputs \autocite{Rasmussen2005}.

\begin{equation}
\bm{f}_\ast \sim \mathcal{N}(\bm{0}, K(\bm{X}_\ast,\bm{X}_\ast)).
\end{equation}

\subsubsection{Prediction with Noise Free Observations}

We are not primarily interested in drawing random functions from the prior, but want to incorporate the knowledge that the training data provides about the function. Let us consider the simple special case where the observations are noise free. That is, we know $\{(\bm{x}_i,f_i)|i=1,\ldots,n\}$. The joint distribution of the training outputs $\bm{f}$ and the test outputs $\bm{f}_\ast$ according to the prior are as follows \autocite{Rasmussen2005}.

\begin{equation}
\begin{bmatrix}
\bm{f} \\
\bm{f}_\ast
\end{bmatrix}
\sim\mathcal{N}\left(\bm{0},
\begin{bmatrix}
K(\bm{X},\bm{X}) & K(\bm{X},\bm{X}_\ast) \\
K(\bm{X}_\ast,\bm{X}) & K(\bm{X}_\ast,\bm{X}_\ast)
\end{bmatrix}
\right).
\end{equation}

If there are $n$ training points and $n_\ast$ test points then $K(\bm{X},\bm{X}_\ast)$ denotes the $n \times n_\ast$ matrix of covariances evaluated at all pairs of training and test points, and similarly for the other entries $K(\bm{X},\bm{X})$, $K(\bm{X}_\ast,\bm{X}_\ast)$, and $K(\bm{X}_\ast,\bm{X})$. To get the posterior distribution over functions we need to restrict this joint prior distribution to contain only those functions which agree with the observed data points \autocite{Rasmussen2005}.

Graphically we may think of generating functions from the prior, and rejecting the ones that disagree with the observations, although this strategy is computationally inefficient. Though in probabilistic terms this operation is simple, corresponding to conditioning the joint Gaussian prior distribution on the observations. In the following, function values $\bm{f}_\ast$ corresponding to test inputs $\bm{X}_\ast$ can be sampled from the joint posterior distribution by evaluating the mean and covariance matrix \autocite{Rasmussen2005}.

\begin{equation}
\begin{split}
\bm{f}_\ast|\bm{X}_\ast,\bm{X},\bm{f}
\sim\mathcal{N}(&K(\bm{X}_\ast,\bm{X})K(\bm{X},\bm{X})^{-1}\bm{f},\\
&K(\bm{X}_\ast,\bm{X}_\ast)-K(\bm{X}_\ast,\bm{X})K(\bm{X},\bm{X})^{-1}K(\bm{X},\bm{X}_\ast)).
\end{split}
\end{equation}

\subsection{Classification}

We have considered regression problems where the targets are real valued. Classification problems aim to assign an input $\bm{x}$ to one of $C$ classes $\mathcal{C}_1,\ldots,\mathcal{C}_C$. These problems can either be binary ($C=2$) or multiclass ($C>2$). Let us focus on probabilistic classification where test predictions take the form of class probabilities. Since generalization to test cases inherently involves some level of uncertainty, it is natural to attempt to make predictions in a way that reflects these uncertainties \autocite{Rasmussen2005}.

Both classification and regression can be viewed as function approximation problems. Unfortunately, the solution of classification problems using Gaussian processes is more demanding than for regression problems. This is because regression problems can assume that the likelihood function is Gaussian. A Gaussian process prior combined with a Gaussian likelihood gives rise to a posterior Gaussian process over functions and everything remains analytically tractable. For classification models where the targets are discrete class labels, the Gaussian likelihood is inappropriate, and we must use methods of approximate inference since exact inference is not feasible \autocite{Rasmussen2005}.

We have seen how Gaussian process regression (GPR) can be obtained by generalizing linear regression. Logistic regression describes an analog of linear regression in the classification case. It is generalized to yield Gaussian process classification (GPC) using again the ideas behind the generalization of linear regression for GPR \autocite{Rasmussen2005}.

For binary discriminative classification the output of a regression model can turn into a class probability using a response function (the inverse of a link function). This squashes its arguments which can lie in the domain $(-\infty,\infty)$ into the domain $[0,1]$ guaranteeing a valid probabilistic interpretation. An example is the linear logistic regression model (\autoref{Appendix:LogReg} for more information) which combines the linear model with the logistic response function \autocite{Rasmussen2005}.

\begin{equation}
p(\mathcal{C}_1|\bm{x})=\lambda(\bm{x}^\top\bm{w}) \quad \text{where} \quad \lambda(z)=\frac{1}{1+\exp(-z)}.
\label{Equation:ResponseFunction}
\end{equation}

Let us consider linear models for binary classification which form the foundation of Gaussian process classification models. We use the labels $y=1$ and $y=-1$ to distinguish the two classes. The likelihood is given by $p(y=1|\bm{x},\bm{w}) = \sigma(\bm{x}^\top\bm{w})$ where $\bm{w}$ is the weights vector and $\sigma(z)$ can be any sigmoid function. When using the logistic $\sigma(z)=\lambda(z)$ from \autoref{Equation:ResponseFunction} the model is simply called logistic regression. A the probability of the two classes must sum to 1, we have $p(y=-1|\bm{x},\bm{w})=1-p(y=1|\bm{x},\bm{w})$. Thus for a data points $(\bm{x}_i,y_i)$ the likelihood is given by $\sigma(\bm{x}_i^\top\bm{w})$ if $y_i=1$, and $1-\sigma(\bm{x}_i^\top\bm{w})$ if $y_i=-1$ \autocite{Rasmussen2005}.

For binary classification the basic idea behind Gaussian process prediction is that we place a GP prior over the latent function $f(\bm{x})$ and then squash this through the logistic function to obtain a prior on $\pi(\bm{x})=p(y=1|\bm{x})=\sigma(f(\bm{x}))$. Note that $\pi$ is a deterministic function of $f$, and since $f$ is stochastic, so if $\pi$. The latent function $f$ is a nuisance function that we do not observe values of itself, and instead only observe the inputs $\bm{X}$ and the class labels $\bm{y}$. We are not particularly interested in the values of $f$, but rather in $\pi$, in particular for test cases $\pi(\bm{x}_\ast)$. The purpose of $f$ is solely to allow a convenient formulation of the model, and the computational goal is to integrate out $f$ \autocite{Rasmussen2005}.

Inference is naturally divided into two steps: First we compute the distribution of the latent variable corresponding to a test case in the following where $p(\bm{f}|\bm{X},\bm{y})=p(\bm{y}|\bm{f})p(\bm{f}|\bm{X})/p(\bm{y}|\bm{X})$ is the posterior over the latent variables \autocite{Rasmussen2005}.

\begin{equation}
p(f_\ast|\bm{X},\bm{y},\bm{x}_\ast) = \int p(f_\ast|\bm{X},\bm{x}_\ast,\bm{f})p(\bm{f}|\bm{X},\bm{y})\mathrm{d}\bm{f}.
\label{Equation:GPClassificationStep1}
\end{equation}

Second we use this distribution over the latent $f_\ast$ to produce a probabilistic prediction \autocite{Rasmussen2005}.

\begin{equation}
\bar{\pi}_\ast = p(y_\ast=1|\bm{X},\bm{y},\bm{x}_\ast) = \int \sigma(f_\ast)p(f_\ast|\bm{X},\bm{y},\bm{x}_\ast)\mathrm{d}f_\ast.
\label{Equation:GPClassificationStep2}
\end{equation}

In the regression case with Gaussian likelihood, computation of predictions was straightforward as the relevant integrals were Gaussian and could be computed analytically. In classification the non-Gaussian likelihood in \autoref{Equation:GPClassificationStep1} makes the integral analytically intractable. Similarly, \autoref{Equation:GPClassificationStep2} can be analytically intractable for certain sigmoid functions. Thus we need to use approximations such as Markov chain Monte Carlo (MCMC) \autocite{Rasmussen2005}.

\subsection{Covariance Functions}

A covariance function encodes our assumptions about the function which we wish to learn. In supervised learning, the notion of similarity between data points is also crucial. It is a basic assumption that points $\bm{x}_p$ and $\bm{x}_q$ that are close are likely to have a similar target value $y$. Under the Gaussian process view it is the covariance function that defines nearness or similarity \autocite{Rasmussen2005}.

A stationary covariance function is a function of $\bm{x}-\bm{x}'$. Thus it is invariant to translations in the input space. If further the covariance function is a function only of $|\bm{x}-\bm{x}'|$ then it is called isotropic, and is thus invariant to all rigid motions. The squared exponential covariance function is isotropic. These functions are also known as radial basis functions (RBFs). Its formulation is as follows where $\ell$ is the characteristic length scale \autocite{Rasmussen2005}.

\begin{equation}
k(\bm{x},\bm{x}') = \exp\left(-\frac{|\bm{x}-\bm{x}'|^2}{2\ell^2}\right).
\end{equation}

Furthermore, if a covariance function depends only on $\bm{x}$ and $\bm{x}'$ through $\bm{x}\cdot\bm{x}'$ it is called a dot product covariance function. The linear covariance function is a special case of the dot product covariance function. The linear covariance function is written as follows where $\bm{c}$ is constant \autocite{Rasmussen2005}.

\begin{equation}
k(\bm{x},\bm{x}') = (\bm{x}-\bm{c})(\bm{x'}-\bm{c}).
\end{equation}

% Squared exponential and linear kernels
\begin{figure}[!htbp]
\centering
\begin{subfigure}{0.49\textwidth}
\includegraphics[width=\linewidth]{appendix/methods/pipelines/assets/gp/exp-quad.png}
\end{subfigure}
\hfill
\begin{subfigure}{0.49\textwidth}
\includegraphics[width=\linewidth]{appendix/methods/pipelines/assets/gp/linear.png}
\end{subfigure}
\caption[Covariance functions]{Example functions drawn from a GP prior with a given covariance function. (Left) The squared exponential covariance function. (Right) The linear covariance function \autocite{Engels2023}.}
\end{figure}

% Additional covariance functions
\begin{figure}[!htbp]
\centering
\begin{subfigure}{0.49\textwidth}
\includegraphics[width=\linewidth]{appendix/methods/pipelines/assets/gp/matern-5-2.png}
\end{subfigure}
\hfill
\begin{subfigure}{0.49\textwidth}
\includegraphics[width=\linewidth]{appendix/methods/pipelines/assets/gp/matern-3-2.png}
\end{subfigure}
\\
\begin{subfigure}{0.49\textwidth}
\includegraphics[width=\linewidth]{appendix/methods/pipelines/assets/gp/exponential.png}
\end{subfigure}
\hfill
\begin{subfigure}{0.49\textwidth}
\includegraphics[width=\linewidth]{appendix/methods/pipelines/assets/gp/polynomial.png}
\end{subfigure}
\caption[Additional covariance functions]{Samples from additional covariance functions. (Top Left) The Mat\'ern $\nicefrac{5}{2}$ kernel. (Top Right) The Mat\'ern $\nicefrac{3}{2}$ kernel. (Bottom Left) The exponential kernel. (Bottom Right) The polynomial kernel \autocite{Engels2023}.}
\end{figure}

A general name for a function $k$ of two arguments mapping a pair of inputs $\bm{x}\in\mathcal{X},\bm{x}'\in\mathcal{X}$ into $\mathbb{R}$ is a kernel. A real kernel is said to be symmetric if $k(\bm{x},\bm{x}')=k(\bm{x}',\bm{x})$. Covariance functions must be symmetric by definition \autocite{Rasmussen2005}.

\subsection{Sparse Approximation}

The best way to perform GP inference and training is to compute a Cholesky decomposition of the $N \times N$ Gram matrix. Unfortunately, this take $\mathcal{O}(N^3)$ time. Instead, we can use an approximation method based on inducing points, also called pseudoinputs, which are like a learned summary of the training data that we can condition on, rather than conditioning on all of it \autocite{Murphy2023}.

Let $\bm{X}$ be the observed inputs, and $\bm{f}_X=f(\bm{X})$ be the unknown vector of function values for which we have observations $\bm{y}$. Let $\bm{f}_\ast$ be the unknown function values at one or more test points $\bm{X}_\ast$. Finally, let us assume we have $M$ additional inputs $\bm{Z}$ with unknown function values $\bm{f}_Z$ also denoted by $\bm{u}$. The exact joint prior has the following form \autocite{Murphy2023}.

\begin{equation}
p(\bm{f}_X,\bm{f}_\ast)
= \int p(\bm{f}_\ast,\bm{f}_X,\bm{f}_Z)\mathrm{d}\bm{f}_Z
= \int p(\bm{f}_\ast,\bm{f}_X|\bm{f}_Z)p(\bm{f}_Z)\mathrm{d}\bm{f}_Z
= \mathcal{N}\left(\bm{0},
\begin{pmatrix}
\bm{K}_{X,X} & \bm{K}_{X,\ast} \\
\bm{K}_{\ast,X} & \bm{K}_{\ast,\ast} \\
\end{pmatrix}
\right).
\end{equation}

We can choose $\bm{f}_Z$ in such a way that it acts as a sufficient statistic for the data, so that we can predict $\bm{f}_\ast$ just using $\bm{f}_Z$ instead of $\bm{f}_X$, and thus we approximate the prior as follows \autocite{Murphy2023}.

\begin{equation}
p(\bm{f}_\ast,\bm{f}_X,\bm{f}_Z)
= p(\bm{f}_\ast|\bm{f}_X,\bm{f}_Z)p(\bm{f}_X|\bm{f}_Z)p(\bm{f}_Z)
\approx p(\bm{f}_\ast|\bm{f}_Z)p(\bm{f}_X|\bm{f}_Z)p(\bm{f}_Z).
\end{equation}

From this we derive the following train and test conditionals \autocite{Murphy2023}.

\begin{align}
p(\bm{f}_X|\bm{f}_Z) &= \mathcal{N}(\bm{f}_X|\bm{K}_{X,Z}\bm{K}_{Z,Z}^{-1}\bm{f}_Z,\bm{K}_{X,X}-\bm{Q}_{X,X}) \\
p(\bm{f}_\ast|\bm{f}_Z) &= \mathcal{N}(\bm{f}_\ast|\bm{K}_{\ast,Z}\bm{K}_{Z,Z}^{-1}\bm{f}_Z,\bm{K}_{\ast,\ast}-\bm{Q}_{\ast,\ast}).
\end{align}

The above equations can be seen as exact inference on noise-free observations $\bm{f}_Z$. To gain computational speedups, we can make further approximations to the terms $\tilde{\bm{Q}}_{X,X}=\bm{K}_{X,X}-\bm{Q}_{X,X}$ and $\tilde{\bm{Q}}_{\ast,\ast}=\bm{K}_{\ast,\ast}-\bm{Q}_{\ast,\ast}$. We can then derive the approximate prior $q(\bm{f}_X,\bm{f}_\ast)=\int q(\bm{f}_X|\bm{f}_Z)q(\bm{f}_\ast|\bm{f}_Z)p(\bm{f}_Z)\mathrm{d}\bm{f}_Z$, which we then condition on the observations in the usual way. All of these approximations result in a training cost of $\mathcal{O}(M^3+NM^2)$, and then take $\mathcal{O}(M)$ time for the predictive mean for each test case, and $\mathcal{O}(M^2)$ time for the predictive variance. Compare this to $\mathcal{O}(N^3)$ training time and $\mathcal{O}(N)$ and $\mathcal{O}(N^2)$ testing time for exact inference \autocite{Murphy2023}.

The deterministic inducing conditional (DIC) approximation, or the subset of regressors (SOR) approximation, results from assuming $\tilde{\bm{Q}}_{X,X}=\bm{0}$ and $\tilde{\bm{Q}}_{\ast,\ast}=\bm{0}$, so the conditionals are deterministic. This can result in an underestimate of the predictive variance. One way to overcome the overconfidence of DIC is to only assume $\tilde{\bm{Q}}_{X,X}=\bm{0}$ but let $\tilde{\bm{Q}}_{\ast,\ast}=\bm{K}_{\ast,\ast}-\bm{Q}_{\ast,\ast}$ be exact. This is called the deterministic training conditional (DTC). Lastly, the fully independent training conditional (FITC) approximation assumes $q(\bm{f}_X|\bm{f}_Z)$ is fully factorized. This throws away less uncertainty than the SOR and DTC methods, since it does not make any deterministic assumptions about the relationship between $\bm{f}_X$ and $\bm{f}_Z$ \autocite{Murphy2023}.

\begin{align}
q_{\mathrm{SOR}}(\bm{f}_X,\bm{f}_\ast) &= \mathcal{N}\left(\bm{0},
\begin{pmatrix}
\bm{Q}_{X,X} & \bm{Q}_{X,\ast} \\
\bm{Q}_{\ast,X} & \bm{Q}_{\ast,\ast} \\
\end{pmatrix}
\right) \\
q_{\mathrm{DTC}}(\bm{f}_X,\bm{f}_\ast) &=
\mathcal{N}\left(\bm{0},
\begin{pmatrix}
\bm{Q}_{X,X} & \bm{Q}_{X,\ast} \\
\bm{Q}_{\ast,X} & \bm{K}_{\ast,\ast} \\
\end{pmatrix}
\right) \\
q_{\mathrm{FITC}}(\bm{f}_X,\bm{f}_\ast) &=
\mathcal{N}\left(\bm{0},
\begin{pmatrix}
\bm{Q}_{X,X}-\operatorname{diag}(\bm{Q}_{X,X}-\bm{K}_{X,X}) & \bm{Q}_{X,\ast} \\
\bm{Q}_{\ast,X} & \bm{K}_{\ast,\ast} \\
\end{pmatrix}
\right).
\end{align}

% Sparse GP comparison
\begin{figure}[!htbp]
\centering
\begin{subfigure}{0.47\textwidth}
\includegraphics[width=\linewidth]{appendix/methods/pipelines/assets/gp/full-gp.png}
\end{subfigure}
\hfill
\begin{subfigure}{0.47\textwidth}
\includegraphics[width=\linewidth]{appendix/methods/pipelines/assets/gp/sparse-dtc.png}
\end{subfigure}
\caption[Sparse Gaussian process]{Sparse GP comparison.}
\end{figure}

The simplest approach to selecting the inducing points from the training set is to throw away some of the data. In this case, we could pick random examples. However, intuitively it makes more sense to try to pick a subset that in some sense covers the original data, so it contains approximately the same information without redundancy. Clustering algorithms such as $k$-means are a popular heuristic approach. We can also use coreset methods which can provably find such an information preserving subset. Furthermore, we can approach learning the inducing points by treating them like kernel hyperparameters, and choosing them so as to maximize the log marginal likelihood \autocite{Murphy2023}.

\end{document}
